\chapter{Transformada de \emph{Fourrier} del campo de \emph{Klein-Gordon}}

	\begin{tikzpicture}
	\fill [left color=red!50, right color=teal!50] (0,0) rectangle (6.5,.1);
	\fill [left color=teal!50, right color=blue!50] (6.5,0) rectangle (11.5,.1);
	\end{tikzpicture}

\vspace{10mm}
\begin{adjustwidth}{50pt}{50pt}
\begin{ejemplo}
\vspace{2mm}
Antes de empezar con la diagonalización del campo de Klein-Gordon, necesitamos ver unos conceptos previos:
\vspace{-3mm}\begin{itemize}
\vspace{-2mm}\item La transformada de Fourrier en el contexto de campos.
\vspace{-2mm}\item Delta de Dirac con transformada de Fourrier
\vspace{-2mm}\item Elemento de volumen en Fourrier que sea invariante Lorentz.	
\end{itemize}

\vspace{2mm}

\end{ejemplo}
\end{adjustwidth}

\vspace{5mm}

\section{Conceptos previos y transformada de Fourrier del campo}



Sea $\{e_1,e_2,\cdots ,e_n\}$ una base ortonormal de $\mathbb R^n$, $e_i\cdot e_j=\delta_{ij}$. Un vector $\vec v$ se expresa como $\vec v =a_1e_1+a_2e_2+\cdots +a_ne_n$

La ecuación e Klein-Gordon en (1+1)-dim es: $\quad \ddot \phi(t,x)+\phi''(t,x)+m^2\phi(t,x)=0$

Una posible solución de esta ecuación, como vimos en el capítulo 33 `Campo escalar con masa', es de la forma $\phi(t,x)=e^{i(\omega t-kx) }$, con la condiciñon de que $\omega^2=k^2+m^2$

Para $k=k_1 \ \to \ \phi_1=e^{i(\omega_1t-k_1x)}$, con $\omega_1=\sqrt{k_1^2+m^2}$ es solución de la ecuación de Klein-Gordon (K-G), y también lo es $\phi_2=e^{i(\omega_2t-k_2x)}$, con $\omega_2=\sqrt{k_2^2+m^2}$, incluso lo es cualquier combinación lineal de ellas, $6e^{i(\omega_1t-k_1x)}-9e^{i(\omega_2t-k_2x)}$ (la ec. K-G es lineal). Luego, $\sum_j b_j\, e^{i(\omega_jt-k_jx)}$ también es solución de la ecuación de Klein-Gordon.

Supongamos que $k$ es una variable continua real y $\omega(k)=\sqrt{k^2+m^2}$. Al tomar $k$ vaores continuos, que no discretos, el sumatorio se convierte en una integral: $\sum \to \int$ y la generalización de la solución general de la ecuación de K-G será:

\begin{equation}
\label{T39SolK-Gctna}
\boldsymbol{
\phi(t,x) \ = \ \int_{-\infty}^{+\infty} \dd k \ b(k)\ e^{i(\omega t-kx)}
}	
\end{equation}

Comprobemos que efectivamente  es una solución de la ec. de K-G \textcolor{gris}{(Derivación e integración definida con límites constantes conmutan)}:

$\displaystyle \ddot \phi=\dv[2]{t} \phi = \int_{-\infty}^{+\infty} \dd k \ b(k) \ (i\omega)^2\ e^{i(\omega t-kx)} \qquad  \qquad \phi''=\dv[2]{x} \phi = \int_{-\infty}^{+\infty} \dd k \ b(k) \ (-ik)^2\ e^{i(\omega t-kx)}$

Sustituyendo en la ecuación de Klein-Gordon,

$\displaystyle 
\int_{-\infty}^{+\infty} \dd x \ b(k)\ (i\omega)^2\ e^{i(\omega t-kx)} -\int_{-\infty}^{+\infty} \dd x \ b(k)\ (ik)^2\ e^{i(\omega t-kx)} + \int_{-\infty}^{+\infty} \dd x \ b(k)\ m^2\  e^{i(\omega t-kx)}= $

$\displaystyle = \int_{-\infty}^{+\infty} \dd x \ b(k)\ [-\omega^2+k^2+m^2]\  e^{i(\omega t-kx)} = 0\, ,  \quad \text{ ya que} \ \ \omega^2=k^2+m^2\, . \ $ Sí verifica le ec. K-G  $ \ \ \Box$


La analogía de la ec. \ref{T39SolK-Gctna} con la de un vector de $\mathbb R^n$ en una base ortonormal nos hace pensar que $e^{i(\omega t-kx)}$ ha de ser, de algún modo, una base ortonormal, habrá de ocurrir algo similar a $e_i \cdot e_j=\delta_{ij}$


$$\displaystyle \vec v \ = \ \textcolor{green}{\sum_{i=1}^n}\,\textcolor{blue}{ a_i}\, \textcolor{red}{e_i} \qquad \quad \sim  \qquad \qquad \phi(t,x) \ = \ \textcolor{green}{\int_{-\infty}^{+\infty} \dd x} \ \textcolor{blue}{b(k)}\ \textcolor{red}{e^{i(\omega t-kx)}}$$

La generalización del producto escalar de vectores ($\vec u \cdot \vec v = \sum_j u_j v_j$) tiene por generalización directa en campos: $\displaystyle \ f(t,x)\, g(t,x)=\int_{-\infty}^{+\infty} f(t,x)\, g(t,x) \, \dd x$

Tomemos dos elementos de la base $e^{i(\omega t-kx)}$ y veamos si es ortonormal. Sean estos $e^{i(\omega_1 t-k_1x)}$ y $e^{i(\omega_2 t-k_2x)}$, su producto escalar será:

$\displaystyle \int_{-\infty}^{+\infty} \dd x  \ e^{i(\omega_1 t-k_1x)} \ e^{i(\omega_2 t-k_2x)} =  \int_{-\infty}^{+\infty} \dd x  \ e^{i[(\omega_1+\omega_2)t-(k_1+k_2)t]} = ( \ $ \textcolor{gris}{$Q=k_1+k_2$} $ \ ) =$

$\displaystyle = e^{i(\omega_1+\omega_2)t} \  \int_{-\infty}^{+\infty} \dd x e^{-iQ\, x} \begin{matrix} \\ = \\ (1) \end{matrix} e^{i(\omega_1+\omega_2)t}\,  2 \pi \, \delta(-Q) \begin{matrix} \\ = \\ (2) \end{matrix}  e^{i(\omega_1+\omega_2)t}\,  2 \pi \, \delta(Q) = e^{i(\omega_1+\omega_2)t}\,  2 \pi \, \delta(k_1+k_2)$

Que, evidentemente no es como esperábamos, no es como un delta de Kronecker.

\begin{footnotesize}
(1) La función delta de Dirac se puede definir como:\footnote{ Cap 13 del curso `Teoría cuántica de campos' de \emph{Javier García}.} $\ \displaystyle \delta(z)=\dfrac 1{2\pi} \int_{-\infty}^{+\infty} \dd m \, e^{i\, m \, z}$

(2) La delta de Dirac es par $\ \delta(-z)=\delta(z)$	
\end{footnotesize}

La solución a este aparente problema está en que cuando se trabaja en el campo complejo el producto escalar no se define como $\int f\,g$ sino como $\int f^* g$. Repitiendo todo el proceso con esta puntualización, se obtiene:
$\ e^{i(\omega_1+\omega_2)t}\, 2\pi\, \delta(-k_1+k_2)$, que sí es correcto \textcolor{gris}{$(k_1=k_2 \to \delta(k_1-k_2)1;\ \ k_1 \neq k_2 \to \delta(k_1-k_2)=0)$}

Para eliminar este molesto factor $\dfrac 1{2 \pi}$, redefinimos el campo como: 

\begin{equation}
\label{T39SolK-Gctna-redef}
\boldsymbol{
\phi(t,x) \ = \ \int_{-\infty}^{+\infty} \dfrac{\dd k}{2\pi} \ b(k)\ e^{i(\omega t-kx)}
}	
\end{equation}


que generalizado a (1+3)-dim, $k\to \vec k=(k_x,k_y,k_z)$ incorporará en el integrando el factor $\dfrac{\dd^3k}{(2\pi)^3}$

Pero la expresión \ref{T39SolK-Gctna-redef} es precisamente la transformada de Fourrier como superposición de ondas planas, pero tal como aparece aquí $\phi \in \mathbb C$ lo cual no puede ser, el campo de Klein-Gordon es real.

Cuando se tiene una magnitud $z\in \mathbb C$ y se quiere trabajar con una real se toma $1/2(z+z^*) \in \mathbb R$, hacíendolo así, definiremos:

\begin{equation}
\label{T39SolK-Gctna-redef-R}
\boldsymbol{
\phi(t,x) \ = \ \dfrac 1 2  \left[\  \int_{-\infty}^{+\infty} \dfrac{\dd k}{2\pi} \ b(k)\ e^{i(\omega t-kx)} \  + \  \int_{-\infty}^{+\infty} \dfrac{\dd k}{2\pi} \ b^*(k)\ e^{-i(\omega t-kx)}  \ \right]
}	
\end{equation}


En el capítulo \ref{campo-escalar-con-masa}, `Campo escalar con masa', vimos el \emph{truco} de De Broglie $p=\hbar k$ y el de Plank $E=\hbar \omega$ que, dado que $p$ y $E$ forman un 4-vector, también lo formarán $k$ y $\omega$. En (3+1)-dim tenemos $\omega t-k_1x-k_2y-k_3z$, como $x^0=ct, x^1=x,x^2=y,x^3=z$ lo podemos escribir como $\omega x^0-k_1x^1-k_2x^2-k_3x^3=k_\mu x^\mu$ y como $k_0=\omega=k^0;\ -k_1=k^1,-k_2=k^2,-k_3=k^3$ se cumple. $\hbar(\omega,\vec k)=(E,\vec p)$ forman un 4-vector. Con esto es seguro que $e^{i(\omega t-kx)}$ queda invariante bajo transformaciones de Lorentz (ya que se trata de un producto escalar $(k_\mu x^\mu)$

Aún nos queda otro problema por resolver, el término $\dfrac {\dd k}{2\pi}$ no es \emph{``manifiestamente''} invariante Lorentz.

Volviendo a dimensión (1+1), si hacemos una transformación de Lorentz,

$\begin{cases} \ x^{0'}=\gamma(x^0-\beta x^1) \\ \ x^{1'}=\gamma(x^1-\beta x^0) \end{cases} \ \gamma=\dfrac 1 {\sqrt{1-\beta^2}}\, ; \ \beta=\dfrac v c \quad \begin{cases} \ k^{0'}=\gamma(k^0-\beta k ^1) \\ \ k^{1'}=\gamma(k^1-\beta k^0) \end{cases} \to \quad \begin{cases} \  \omega'=\gamma(\omega-\beta k^1) \\ \ k^{1'}=\gamma(k^1-\beta k^0) \end{cases}$

De la última ecuación, $\ \displaystyle \dd {K'}=\gamma (\dd k - \beta \dd \omega) \neq \dd k \ \ $ luego $\ \ \dd k' \begin{matrix} \text{\tiny{Trans. Lorentz}} \\ \cancel{\longrightarrow} \\ \\ \end{matrix}  \dd k\, , \ $ no es `manifiestamente' invariante Lorentz.

Para resolver este contratiempo y modificar el factor $\dfrac {\dd k}{2\pi}$ para que sí sea `manifiestamente' invariante Lorentz, lo que haremos  es lo siguiente:

$\displaystyle \dv{K'}{k}=\gamma \left( 1-\beta \dv{\omega}{k} \right) = \ \textcolor{gris}{(\omega=\sqrt{k^2+m^2})} \ = \gamma \left( 1-\beta \dfrac{\cancel{2}k}{\cancel{2}\sqrt{k^2+m^2}} \right) =\gamma \left( 1-\beta \dfrac{k}{\omega} \right) =\dfrac{\gamma}{\omega} (\omega-\beta k) =\dfrac{\omega'}{\omega}$

Despejando, $\quad \dfrac{\dd k'}{\omega'} \ = \ \dfrac{\dd k}{\omega} \ $ por lo que $\ \dfrac{\dd k}{\omega} $ en `manifiestamente' invariante ante una transformación de Lorentz. Lo que debemos hacer es volver a redefinir el campo como:

\begin{equation}
\label{T39SolK-Gctna-redef-R-IL}
\boldsymbol{
\phi(t,x) \ = \ \dfrac 1 2  \left[\  \int_{-\infty}^{+\infty} \dfrac{\dd k}{2\pi \omega} \ \omega \  b(k)\ e^{i(\omega t-kx)} \  + \  \int_{-\infty}^{+\infty} \dfrac{\dd k}{2\pi \omega} \  \omega \ b^*(k)\ e^{-i(\omega t-kx)}  \ \right]
}	
\end{equation}

Al ser $b=b(k)$ y también $\omega=\omega(k)$, llamamos $\omega \ b(k) \ = \ \sqrt{k^2+m*2}\ b(k) \ = \ \boldsymbol{a(k)} \ $ y problema solucionado, el campo será:

\vspace{3mm}
\begin{equation}
\label{T39SolK-Gctna-redef-R-a}
\boldsymbol{ \boxed{ \  \subrayado{
\phi(t,x) \ = \ \dfrac 1 2   \int_{-\infty}^{+\infty} 
\dfrac{\dd k}{2\pi \omega} \ \left[ \ 
\ a(k)\ e^{i(\omega t-kx)} \  + \ a^*(k)\ e^{-i(\omega t-kx)}
\ \right]
} \ } }	
\end{equation}
\vspace{3mm}

Como vamos a necesitar $\dot \phi = \pi$, derivemos el campo respecto el tiempo:

$\displaystyle \pi(t,x)\ = \ \dv{\phi}{t} \ = \ 
\dfrac 1 2   \int_{-\infty}^{+\infty} 
\dfrac{\dd k}{2\pi \omega} \ \left[ \ 
\ a(k)\ (i\omega)\ e^{i(\omega t-kx)} \  + \ a^*(k)\ (-i\omega)\ e^{-i(\omega t-kx)}
\ \right]
$

\vspace{3mm}
\begin{equation}
\label{T39Pi}
\boldsymbol{ \boxed{ \ \subrayado{
 \pi(t,x)\ = \
\dfrac 1 2   \int_{-\infty}^{+\infty} 
\dfrac{\dd k}{2\pi \omega} \ \ (i\omega) \ \left[ \ 
\ a(k)\ e^{i(\omega t-kx)} \ -  \ a^*(k)\ e^{-i(\omega t-kx)}
\ \right]
} \ } }		
\end{equation}
\vspace{3mm}

Aún estamos en el `mundo clásico' (no cuántico) $\ \leftrightarrow \ a(k)\, , a^+(k)\ $ son números (en Q serán operadores).


\vspace{10mm}

\section{Ejercicio propuesto}
\vspace{5mm}
\begin{ejercicio}
\begin{adjustwidth}{30pt}{20pt}
\vspace{2mm}
Demostrar que $\ \dfrac{\dd^3 K}{\omega}\ $ 	 es, manifiestamente, un invariante Lorentz, es decir, que 

$\ \displaystyle \dfrac{\dd^3 K}{\omega}\  = \ \dfrac{\dd^3 K'}{\omega'}\, , \ $ siendo $\ k=\sqrt{k_x^2+k_y^2+k_z^2+m^2}$
\vspace{2mm}
\end{adjustwidth}
\end{ejercicio}

\vspace{5mm}
\color{MidnightBlue}
4-vector: $\ k^{\mu}=(\omega, k^1,k^2,k^3) \ \to \text{ TL: } \ \ k^{\mu'}=(\omega', (k^1)',(k^2)',(k^3)')$

$\begin{cases} 
\ \omega' = \gamma (\omega - \beta k^1) \\
\ (k^1)'=\gamma (-\beta \omega + k^1) \\
\ (k^2)' = k^2 & \to \dd (k^2)' = \dd k^2	\\
\ (k^3)' = k^3 & \to \dd (k^3)' = \dd k^3
\end{cases} \qquad \gamma = \dfrac 1{\sqrt{1-\beta^2}} \quad \beta=\dfrac v c$

$\dfrac{\dd^3 k}{\omega}= \dfrac{\dd k^1\, \dd k^2\, \dd k_3}{\omega}$


$ \left. \begin{matrix}
\displaystyle \dv{(k^1)'}{k^1}=\gamma \left(-\beta \dv{\omega}{k^1} + 1 \right)
\qquad \qquad \qquad \qquad   \\
\displaystyle \dv{\omega}{k^1}= \dfrac{\cancel{2}k^1}{\cancel{2} \sqrt{(k^1)^2+(k^2)^2+(k^3)^2+m^2}}=\dfrac{k^1}{\omega} 
\end{matrix} \right\} \ \Rightarrow \ \dv{(k^1)'}{k^1}=\gamma \left( 1-\beta \dfrac {k'}{\omega} \right) = \dfrac 1 \omega [\gamma (\omega - \beta k^1)] = \dfrac 1 \omega \, \omega' $

Luego $\ \displaystyle \boldsymbol{ \dfrac{\dd k'}{\omega'} = \dfrac{\dd k}{\omega}}$, como $\ \boldsymbol{ \dd (k^2)'=\dd k^2 \text{ y } \dd (k^3)'=\dd k^3}\ $ entonces,

$\boldsymbol{ \dfrac{\dd^3 k}{\omega}=} \dfrac{\dd k^1\, \dd k^2\, \dd k_3}{\omega} = \dfrac{\dd k^1}{\omega} \, \dd k^2\, \dd k_3 = \dfrac{\dd (k^1)'}{\omega'} \, \dd( k^2)' \, \dd (k^3)' = \boldsymbol{ \dfrac{\dd^3 k'}{\omega'}}$ \hspace{5cm} $\Box$






























\color{black}

